\documentclass[12pt]{article}
\usepackage{amsmath,amssymb,amsfonts}
\usepackage[margin=1in]{geometry}
\usepackage{bm}

\newcommand{\vect}[1]{\mathbf{#1}}
\newcommand{\ehat}[1]{\hat{\vect{e}}_{#1}}

\begin{document}

\section*{Problem 16}

\textbf{Problem:} A charge $q$ is uniformly distributed upon a circle that has radius $R$, is concentric with the $z$-axis, and lies in the plane $z = 0$. A charge $-q$ is located at the center of the circle. Find:
\begin{itemize}
\item[(i)] the electric monopole moment,
\item[(ii)] the electric dipole moment vector, and
\item[(iii)] the electric quadrupole moment tensor.
\end{itemize}
Using the symmetry of the quadrupole tensor and the fact that its trace vanishes, as well as the inherent symmetry of the physical system under consideration, one can reduce the labor enormously.

\begin{itemize}
\item[(b)] Determine for this charge distribution the multipole expansion for the potential at distance $r$ from the origin for which $r \gg R$; carry the expansion through the quadrupole contribution, and express the spatial dependence in the final result in terms of the spherical coordinates $r$, $\theta$, $\psi$.
\end{itemize}

\bigskip
\textbf{Solution:}

\subsection*{Setup}

The charge distribution consists of:
\begin{enumerate}
\item A charge $-q$ at the origin.
\item A charge $+q$ uniformly distributed on a circle of radius $R$ in the $z = 0$ plane, centered at the origin.
\end{enumerate}

For the ring, the linear charge density is $\lambda = q/(2\pi R)$. A point on the ring is at position $\vect{r}' = R\cos\phi'\,\ehat{x} + R\sin\phi'\,\ehat{y}$, with line element $dl = R\,d\phi'$.

\subsection*{(i) Monopole moment}

The monopole moment (total charge) is:
\[
Q = q + (-q) = \boxed{0}.
\]

\subsection*{(ii) Dipole moment}

The dipole moment is (Table 1.2, Eq.~1.103):
\[
\vect{p} = \int \vect{r}'\,\rho(\vect{r}')\,d\tau'.
\]

For the point charge at the origin: contribution $= (-q)\cdot\vect{0} = \vect{0}$.

For the ring:
\[
\vect{p}_{\text{ring}} = \int_0^{2\pi} (R\cos\phi'\,\ehat{x} + R\sin\phi'\,\ehat{y})\,\frac{q}{2\pi R}\,R\,d\phi' = \frac{qR}{2\pi}\int_0^{2\pi}(\cos\phi'\,\ehat{x} + \sin\phi'\,\ehat{y})\,d\phi'.
\]

Both integrals vanish: $\int_0^{2\pi}\cos\phi'\,d\phi' = \int_0^{2\pi}\sin\phi'\,d\phi' = 0$.

Therefore:
\[
\boxed{\vect{p} = \vect{0}}.
\]

This is expected from the azimuthal symmetry.

\subsection*{(iii) Quadrupole moment tensor}

The quadrupole moment tensor is defined as (Table 1.2):
\[
Q_{ij} = \int (3x_i'x_j' - r'^2\delta_{ij})\,\rho(\vect{r}')\,d\tau'.
\]

The point charge at the origin contributes zero (since $\vect{r}' = 0$).

For the ring, $r'^2 = R^2$, $x_1' = R\cos\phi'$, $x_2' = R\sin\phi'$, $x_3' = 0$. We need the following angular averages over the ring:
\begin{align}
\langle\cos^2\phi'\rangle &= \frac{1}{2\pi}\int_0^{2\pi}\cos^2\phi'\,d\phi' = \frac{1}{2}, \\
\langle\sin^2\phi'\rangle &= \frac{1}{2}, \\
\langle\cos\phi'\sin\phi'\rangle &= 0.
\end{align}

\textbf{Using symmetry:} The system has azimuthal symmetry about the $z$-axis, so $Q_{ij}$ must be diagonal with $Q_{11} = Q_{22}$ (the $x$ and $y$ directions are equivalent). Also, the trace must vanish: $Q_{11} + Q_{22} + Q_{33} = 0$. So we only need one independent component.

Compute $Q_{33}$:
\begin{align}
Q_{33} &= \int (3x_3'^2 - r'^2)\,\rho\,d\tau'.
\end{align}

On the ring, $x_3' = 0$ and $r'^2 = R^2$, so:
\[
Q_{33} = q\,(3\cdot 0 - R^2) = -qR^2.
\]

From the tracelessness condition and $Q_{11} = Q_{22}$:
\[
2Q_{11} + Q_{33} = 0 \quad \Longrightarrow \quad Q_{11} = Q_{22} = \frac{qR^2}{2}.
\]

All off-diagonal elements vanish by the azimuthal symmetry. Therefore:
\[
\boxed{Q_{ij} = \frac{qR^2}{2}\begin{pmatrix} 1 & 0 & 0 \\ 0 & 1 & 0 \\ 0 & 0 & -2 \end{pmatrix}}.
\]

\subsection*{(b) Multipole expansion of the potential}

The general multipole expansion (Eq.~1.103) is:
\[
\varphi(\vect{r}) = \frac{Q}{r} + \frac{\vect{p}\cdot\hat{\vect{r}}}{r^2} + \frac{1}{6}\sum_{ij}\frac{Q_{ij}\,\hat{r}_i\hat{r}_j}{r^3} + \cdots
\]

Since $Q = 0$ and $\vect{p} = \vect{0}$, the first nonvanishing term is the quadrupole:
\[
\varphi(\vect{r}) = \frac{1}{6}\frac{Q_{ij}\,\hat{r}_i\hat{r}_j}{r^3}.
\]

Compute $Q_{ij}\hat{r}_i\hat{r}_j$:
\begin{align}
Q_{ij}\hat{r}_i\hat{r}_j &= Q_{11}\hat{r}_1^2 + Q_{22}\hat{r}_2^2 + Q_{33}\hat{r}_3^2 \\
&= \frac{qR^2}{2}(\hat{r}_1^2 + \hat{r}_2^2) - qR^2\hat{r}_3^2 \\
&= \frac{qR^2}{2}(1 - \hat{r}_3^2) - qR^2\hat{r}_3^2 \\
&= \frac{qR^2}{2} - \frac{qR^2}{2}\hat{r}_3^2 - qR^2\hat{r}_3^2 \\
&= \frac{qR^2}{2} - \frac{3qR^2}{2}\hat{r}_3^2 \\
&= \frac{qR^2}{2}(1 - 3\cos^2\theta),
\end{align}
where $\hat{r}_3 = \cos\theta$.

Therefore:
\[
\varphi(r,\theta) = \frac{1}{6}\cdot\frac{qR^2(1 - 3\cos^2\theta)}{2r^3} = \frac{qR^2(1 - 3\cos^2\theta)}{12\,r^3}.
\]

Using the Legendre polynomial $P_2(\cos\theta) = \frac{1}{2}(3\cos^2\theta - 1)$:
\[
\boxed{\varphi(r,\theta) = -\frac{qR^2}{6r^3}\,P_2(\cos\theta)}.
\]

The potential has no angular dependence on $\psi$ (the azimuthal angle), as expected from the axial symmetry. The potential falls off as $1/r^3$, characteristic of a pure quadrupole. \hfill $\blacksquare$

\end{document}
